\documentclass[10pt,twocolumn]{article}
\usepackage[T1]{fontenc}
\usepackage[utf8]{inputenc}
\usepackage{lmodern}
\usepackage[margin=1.8cm,top=2.4cm,bottom=2cm,columnsep=0.6cm]{geometry}
\usepackage{fancyhdr}
\usepackage{titlesec}
\usepackage{booktabs}
\usepackage{amsmath,amssymb,amsfonts}
\usepackage{microtype}
\usepackage{xcolor}
\usepackage{mdframed}
\usepackage{parskip}
\usepackage{enumitem}
\usepackage{hyperref}

\definecolor{rulered}{RGB}{180,30,30}
\definecolor{boxgray}{RGB}{245,245,245}
\definecolor{keywordgray}{RGB}{100,100,100}

\pagestyle{fancy}
\fancyhf{}
\fancyhead[L]{\textsc{\small Claude Mag}}
\fancyhead[C]{\small\textit{Machine Learning Research Digest}}
\fancyhead[R]{\small 2026-07-20}
\fancyfoot[C]{\small\thepage}
\renewcommand{\headrulewidth}{0.8pt}

\titleformat{\section}{\normalsize\bfseries\color{rulered}}{}{0pt}{}%
  [\vspace{-4pt}\color{rulered}\rule{\columnwidth}{0.4pt}]
\titlespacing{\section}{0pt}{8pt}{4pt}

\hypersetup{colorlinks=true,urlcolor=rulered,linkcolor=black}

\begin{document}
\setlength{\parindent}{0pt}
\setlength{\parskip}{4pt}

{\small\color{keywordgray}\textbf{Keywords:}
  diffusion language models \textbullet{}
  policy gradient \textbullet{}
  masked diffusion \textbullet{}
  RLHF}\par\vspace{4pt}

{\LARGE\bfseries\raggedright
  Mask-Aware Policy Gradients}\par\vspace{4pt}

{\large\itshape\raggedright
  UT Austin Corrects a Fundamental Bias in RL Training for
  Masked Diffusion Language Models by Properly Accounting
  for Both Token and Masking Decisions}\par\vspace{6pt}

{\small\textbf{By} Haran Raajesh, Kulin Shah, Adam Klivans,
  Philipp Kr\"{a}henb\"{u}hl (University of Texas at
  Austin)}\par\vspace{2pt}

{\small\color{keywordgray}arXiv:2607.15200 \textbullet{}
  COLM 2026 \textbullet{} July 17, 2026}
\par\vspace{2pt}\noindent{\color{rulered}\rule{\columnwidth}{0.6pt}}\vspace{6pt}

Masked Diffusion Language Models (MDLMs) generate text by iteratively
unmasking tokens over multiple steps. Applying reinforcement learning
to MDLMs --- to improve performance on math reasoning, code generation,
or instruction following --- requires a valid policy gradient. But all
existing RL approaches for MDLMs compute a policy gradient that ignores
one of the model's two fundamental decisions: not just which tokens
to predict, but which positions to remask at each step. TRACE-MDLM
corrects this omission with a two-stage action MDP formulation, raising
GSM8K accuracy from 79.3\% to 87.1\% and MBPP from 41.2\% to 53.4\%.

\begin{mdframed}[backgroundcolor=boxgray,linecolor=rulered,linewidth=1pt,
    innerleftmargin=6pt,innerrightmargin=6pt,
    innertopmargin=6pt,innerbottommargin=6pt]
\textbf{\small AT A GLANCE}\par\vspace{4pt}
\begin{description}[leftmargin=0pt,labelsep=4pt,itemsep=2pt,parsep=0pt,
    font=\small\bfseries]
  \item[Problem:] RL for MDLMs ignores the masking schedule decision
    in the policy gradient, producing a biased update that
    misattributes credit to token choices.
  \item[Why it matters:] MDLMs are an increasingly competitive
    alternative to autoregressive LMs; reliable RL is necessary
    for closing the remaining gap to AR models on reasoning tasks.
  \item[Method:] Two-stage action MDP: the policy gradient decomposes
    into a token term and a masking term; a lightweight MLP head
    is trained to parameterize the masking policy.
  \item[Result:] 87.1\% on GSM8K (+7.8 pp over token-only baseline);
    53.4\% on MBPP (+12.2 pp); 40\% lower gradient variance.
  \item[Evidence:] Evaluated on GSM8K and MBPP; ablations on 3B-scale
    MDLM; accepted at COLM 2026.
\end{description}
\end{mdframed}

\section{The Big Picture}

Masked Diffusion Language Models are a promising alternative to
autoregressive (AR) LMs: they are intrinsically parallelizable
across token positions, support flexible bidirectional context, and
can be sampled at adjustable quality-speed tradeoffs by varying the
number of denoising steps. MDLMs have closed much of the gap to
AR models on language modeling benchmarks, but lag behind on
structured reasoning tasks like mathematics and code.

Reinforcement learning from verifiable rewards --- the technique
that drove major gains for AR reasoning models via GRPO, PPO, and
similar methods --- is the natural next step for MDLMs. The obstacle
is that naive application of AR policy gradient methods to MDLMs
produces a biased estimator, because it ignores a decision that
is structurally part of the MDLM generation process: the masking
schedule.

\section{Why Existing Approaches Fall Short}

Standard RL for AR models computes the log-probability of each
generated token and uses these log-probabilities directly in the
policy gradient. For MDLMs, the analogous computation scores each
predicted token at each unmasked position --- ignoring that the
model also decided which positions to keep unmasked (and which to
remask) at each step.

The masking schedule is not a fixed algorithmic artifact: it is
a learned policy decision. A given sequence of $T$ denoising steps
involves $T$ masking decisions, each dependent on the model's output
and the task. Ignoring these decisions is equivalent to computing
the policy gradient for an AR model while ignoring every other
token in the generation --- the gradient is biased and the resulting
training signal miscredits token choices for outcomes that were
actually driven by masking decisions.

Prior approaches including SPG and GRPO-for-MDLMs apply variance
reduction techniques to the token gradient without addressing this
structural bias.

\section{Core Mechanism}

The paper formalizes MDLM generation as a two-stage MDP per
denoising step. At step $t$, the model observes a partially masked
sequence $\mathbf{x}_t$ and takes two sequential actions:

\begin{enumerate}[itemsep=2pt,parsep=0pt]
  \item \textbf{Token action $a^{\text{tok}}_t$:} For each masked
    position $i \in \mathcal{M}_t$, predict the unmasked token
    $x_i$ from the MDLM network. This recovers a fully unmasked
    intermediate sequence $\hat{\mathbf{x}}_t$.
  \item \textbf{Masking action $a^{\text{mask}}_t$:} Select a
    subset of positions in $\hat{\mathbf{x}}_t$ to remask,
    producing $\mathbf{x}_{t+1}$ for the next step. This is
    parameterized by a lightweight MLP over the MDLM hidden
    states.
\end{enumerate}

The full policy gradient decomposes as:
\[
\nabla_\theta J = \mathbb{E}_\pi \left[ \hat{A}_t \!\left(
  \nabla_\theta \log \pi_\theta(a^{\text{tok}}_t \mid \mathbf{x}_t)
  + \nabla_\theta \log \pi_\theta(a^{\text{mask}}_t \mid \hat{\mathbf{x}}_t)
\right) \right]
\]
where $\hat{A}_t$ is the advantage estimate at step $t$.

\section{Technical Formulation}

The masking policy $\pi_\theta(a^{\text{mask}}_t \mid \hat{\mathbf{x}}_t)$
is a factored Bernoulli over positions:
\[
\pi_\theta(a^{\text{mask}}_t \mid \hat{\mathbf{x}}_t)
= \prod_{i=1}^{L} \text{Bernoulli}\!\left(m_i;\; \sigma(\mathbf{W}
  \mathbf{h}_i + b)\right)
\]
where $\mathbf{h}_i$ is the MDLM hidden state at position $i$ after
the token action, $\mathbf{W} \in \mathbb{R}^{1 \times d}$ and
$b \in \mathbb{R}$ are the MLP parameters, and $m_i \in \{0, 1\}$
is the masking decision. The MLP is trained jointly with the
denoising network.

The advantage $\hat{A}_t$ is estimated using Generalized Advantage
Estimation (GAE) with $\lambda = 0.95$ and a value function head
trained on the MDLM encoder representations.

\section{Experimental Findings}

\begin{center}
\small
\begin{tabular}{lcc}
\toprule
Method & GSM8K & MBPP \\
\midrule
MDLM-3B (no RL) & 52.1\% & 28.7\% \\
Token-only REINFORCE & 79.3\% & 41.2\% \\
SPG & 81.4\% & 44.7\% \\
\textbf{Mask-Aware PG (ours)} & \textbf{87.1\%} & \textbf{53.4\%} \\
\bottomrule
\end{tabular}
\end{center}

The masking term alone contributes approximately 3--4 percentage
points; jointly optimizing both token and masking gradients
provides the full 7--12 percentage point improvement over
token-only methods.

Gradient norm variance is 40\% lower with Mask-Aware PG,
explaining the smoother training curves and reduced need for
gradient clipping.

\section{Ablations and Interpretation}

\textbf{Token gradient only:} Removing the masking term
reproduces the SPG/REINFORCE baseline, confirming the masking
term as the source of improvement.

\textbf{Masking term only:} Training with only the masking
gradient (no token gradient) fails to converge, confirming
both terms are jointly necessary.

\textbf{Masking schedule analysis:} Models trained with
Mask-Aware PG learn to unmask tokens in a semantically
coherent order: in math problems, nouns and verbs are
unmasked first, modifiers last. This learned curriculum
differs qualitatively from the uniform-random schedule used
at initialization.

\textbf{Scaling:} The advantage grows with problem difficulty.
For easy single-step GSM8K problems the gap is $<$3 percentage
points; for multi-step problems requiring 4+ intermediate
computations the gap is 8+ percentage points, consistent
with the masking schedule being most important when sequential
reasoning structure matters.

\section*{Reference}

Haran Raajesh, Kulin Shah, Adam Klivans, Philipp Kr\"{a}henb\"{u}hl.
\textbf{Mask-Aware Policy Gradients for Diffusion Language Models}.
COLM 2026 / arXiv:2607.15200, July 2026.
\url{https://arxiv.org/abs/2607.15200}

\end{document}
